\chapter{Tablas completas de resultados}
\label{ap:tablas_resultados}

Este apéndice recoge tablas de resultados ampliadas que sirven de soporte a la evaluación del Capítulo~\ref{cap:experimentos}. Algunas consolidan en un mismo lugar cifras ya presentadas allí y otras añaden detalle, en concreto los recuentos de confusión completos de la detección binaria y el desglose por clase de la línea base de aprendizaje automático. El análisis y la interpretación de estos números, con la discusión de los aciertos y los errores, se desarrollan en el Capítulo~\ref{cap:experimentos}.

Todas las cifras provienen de los ficheros de resultados del proyecto. El escaneo SSH (\texttt{anomaly-sshscan}) apenas tiene muestras etiquetadas en \texttt{august.week1}, por lo que su evaluación representativa se hace en \texttt{april.week2}. La clase de campaña SMTP (\texttt{anomaly-spam}) se descartó del análisis y no se incluye como resultado activo.

Las tablas de este apéndice proceden de ficheros CSV derivados de los resultados del clasificador y de la línea base. Son de tres tipos. El primero es un \textbf{resumen binario}, con los recuentos de confusión (VP, FP, FN y VN) y las métricas de precisión, recall y $F_1$. El segundo es un \textbf{resumen por familia y subtipo}, con la clase, las trazas etiquetadas, las predichas, los aciertos y sus métricas. El tercero es un \textbf{resumen de la línea base de aprendizaje automático}, con el modelo, la clase, la precisión, el recall, el $F_1$ y el soporte. Todas estas tablas son de elaboración propia a partir de los resultados experimentales del proyecto, que forman parte del material complementario.

\section{Evolución de las versiones del clasificador}
\label{sec:apc_binaria}

La Tabla~\ref{tab:apc_binaria} resume la evolución de las cinco versiones del clasificador. Para cada una indica el cambio principal y, en las que disponen de un resumen comparable, su detección binaria (precisión, recall y $F_1$) sobre el conjunto principal de estudio (\texttt{august.week1}). De la v2 a la v5 la evaluación se hizo con la misma metodología a nivel de flujo. La v1 fue una primera versión exploratoria que no dejó un resumen binario homogéneo, por lo que no se reportan sus métricas.

\begin{table}[!htbp]
\centering
\small
\begin{tabular}{|l|p{4.4cm}|c|c|c|p{3.6cm}|}
\hline
\textbf{Versión} & \textbf{Cambio principal} & \textbf{P} & \textbf{R} & \textbf{F1} & \textbf{Observación o limitación} \\
\hline
v1 & Reglas locales (ventana de $\pm$30 filas) con señales básicas. & --- & --- & --- & Etapa exploratoria, sin un resumen binario homogéneo comparable. \\
\hline
v2 & Decisión jerárquica (ataque/background, familia y subtipo). & 0,961 & 0,991 & 0,976$^{*}$ & Buen binario, pero el subtipo \texttt{scan11}/\texttt{scan44} muy inestable. \\
\hline
v3 (base estable) & Segundo pase global por ventana. & 0,930 & 0,991 & 0,960$^{*}$ & Estabiliza el escaneo vertical; aún no detecta el escaneo SSH. \\
\hline
v4 & Reparto de confianza alto/bajo en la botnet. & 0,930 & 0,991 & 0,960$^{*}$ & Experimental; binario igual que la v3, mejora no concluyente. \\
\hline
v5 (final) & Tercer pase global por origen (fan-out SSH al puerto 22). & 0,926 & 0,991 & 0,957 & Recupera el escaneo SSH; añade falsos positivos en semanas sin SSH. \\
\hline
\end{tabular}
\caption{Evolución de las versiones del clasificador: cambio principal y detección binaria (ataque/normal) en \texttt{august.week1}. $^{*}$El F1 marcado se deriva de la precisión y el recall, ya que esos resúmenes no lo almacenan. La v1 fue una primera versión exploratoria sin un resumen binario homogéneo.}
\label{tab:apc_binaria}
\end{table}

Entre la base estable v3 y la versión final v5, la diferencia está en los falsos positivos: la v5 añade 8\,337 falsas alarmas respecto a la v3 (de 137\,365 a 145\,702), atribuibles al tercer pase de escaneo SSH en una semana donde ese ataque apenas aparece etiquetado. Los recuentos de confusión completos (VP, FP, FN y VN) de la v3 y la v5 se recogen en el Capítulo~\ref{cap:experimentos}.

\section{Resultados por subtipo de la v5}
\label{sec:apc_familia}

La Tabla~\ref{tab:apc_familia} recoge, por subtipo, las trazas etiquetadas, las predichas y los aciertos sobre \texttt{august.week1}, junto con las métricas derivadas.

\begin{table}[ht]
\centering
\small
\begin{tabular}{|l|r|r|r|c|c|c|}
\hline
\textbf{Subtipo} & \textbf{Etiquetadas} & \textbf{Predichas} & \textbf{Aciertos} & \textbf{P} & \textbf{R} & \textbf{F1} \\
\hline
anomaly-udpscan & 1\,001\,413 & 1\,067\,045 & 1\,001\,219 & 0,938 & 1,000 & 0,968 \\
\hline
scan11 & 102\,572 & 133\,989 & 102\,220 & 0,763 & 0,997 & 0,864 \\
\hline
scan44 & 382\,351 & 399\,411 & 304\,396 & 0,762 & 0,796 & 0,779 \\
\hline
dos & 247\,230 & 217\,724 & 120\,635 & 0,554 & 0,488 & 0,519 \\
\hline
nerisbotnet & 99\,762 & 16\,462 & 4\,420 & 0,269 & 0,044 & 0,076 \\
\hline
anomaly-sshscan & 44 & 8\,397 & 0 & 0,000 & 0,000 & 0,000 \\
\hline
\end{tabular}
\caption{Resultados por subtipo de la v5 en \texttt{august.week1}. El escaneo SSH no es representativo en esta semana, con solo 44 trazas etiquetadas, y se evalúa en \texttt{april.week2} (Tabla~\ref{tab:apc_general}).}
\label{tab:apc_familia}
\end{table}

\section{Generalización por semana}
\label{sec:apc_general}

La Tabla~\ref{tab:apc_general} resume la detección binaria de la v5 en las tres semanas evaluadas. Son las cifras del resumen utilizado en el Capítulo~\ref{cap:experimentos}. No existe un fichero independiente de generalización en el repositorio, por lo que estos valores son los mismos que se presentan e interpretan en ese capítulo.

\begin{table}[ht]
\centering
\small
\begin{tabular}{|l|c|c|c|p{5.0cm}|}
\hline
\textbf{Conjunto} & \textbf{P} & \textbf{R} & \textbf{F1} & \textbf{Familias presentes} \\
\hline
\texttt{august.week1} & 0,926 & 0,991 & 0,957 & Núcleo (udpscan, scan11/scan44, dos, nerisbotnet). \\
\hline
\texttt{august.week2} & 0,872 & 0,747 & 0,805 & Escaneo vertical, dos, nerisbotnet. Sin udpscan. \\
\hline
\texttt{april.week2} & 0,845 & 0,903 & 0,873 & Sobre todo escaneo SSH. \\
\hline
\end{tabular}
\caption{Detección binaria de la v5 por semana. Cifras del resumen usado en el Capítulo~\ref{cap:experimentos}.}
\label{tab:apc_general}
\end{table}

\section{Línea base de aprendizaje automático por clase}
\label{sec:apc_ml}

La Tabla~\ref{tab:apc_ml} desglosa por clase activa el F1 de cada modelo de la línea base, que el Capítulo~\ref{cap:experimentos} resume solo como F1 macro. El cálculo se hace sobre una muestra equilibrada con el mismo número de ejemplos por clase.

\begin{table}[ht]
\centering
\resizebox{\textwidth}{!}{%
\begin{tabular}{|l|c|c|c|c|c|c|c|}
\hline
\textbf{Modelo} & \textbf{background} & \textbf{anomaly-udpscan} & \textbf{dos} & \textbf{nerisbotnet} & \textbf{scan11} & \textbf{scan44} & \textbf{F1 macro} \\
\hline
Regresión logística & 0,703 & 0,982 & 0,995 & 0,869 & 0,705 & 0,304 & 0,760 \\
\hline
KNN & 0,821 & 0,990 & 0,999 & 0,987 & 0,934 & 0,924 & 0,942 \\
\hline
SVM & 0,767 & 0,990 & 0,995 & 0,912 & 0,707 & 0,311 & 0,780 \\
\hline
MLP & 0,833 & 0,990 & 1,000 & 0,992 & 0,801 & 0,721 & 0,889 \\
\hline
\end{tabular}}
\caption{F1 por clase activa de la línea base de aprendizaje automático (muestra equilibrada, 2\,778 ejemplos por clase). Recoge los cuatro modelos de la comparación principal de la memoria.}
\label{tab:apc_ml}
\end{table}

Las clases activas son el tráfico de fondo (\emph{background}) y los cinco tipos de ataque presentes en la muestra. Quedan fuera \texttt{anomaly-spam} y la etiqueta de listas negras (\emph{blacklist}), coherentemente con el enfoque del trabajo.

Estas tablas son el soporte numérico ampliado de la evaluación.
